\subsection{Conventional supervision and engineered customer profiles}

Earlier non-technical-loss research established the customer as a useful prediction unit. Messinis and Hatziargyriou reviewed detection families, data sources, and evaluation issues, showing that non-technical-loss studies differ substantially in assumptions and protocol \cite{Messinis2018NTLReview}. Buzau et al. used utility smart-meter data and inspected non-technical-loss targets with supervised learning and engineered features, demonstrating the practical strength of customer-profile models \cite{Buzau2019NTL}. Shehzad et al. combined big-data processing and genetic-algorithm optimization for supervised electricity-theft classification on SGCC histories \cite{Shehzad2022GeneticETD}. These approaches motivate the FULL\_48 comparator in AMP: a strong engineered customer profile should not be displaced by a learned method unless the comparison is governed and informative.

\subsection{Deep temporal representation learning}

Deep models increasingly transform a long customer history before classification. Zheng et al. combined wide and deep convolutional views of SGCC consumption, including one- and two-dimensional temporal representations \cite{Zheng2018WideDeep}. Hasan et al. combined convolutional feature extraction with an LSTM and used SMOTE to address imbalance \cite{Hasan2019CNNLSTM}. More broadly, the time-series classification review by Ismail Fawaz et al. documents convolutional, recurrent, and hybrid mechanisms for learning temporal representations \cite{Fawaz2019TimeSeriesReview}. This body of work supports learning from temporal structure, but it does not by itself resolve the label-granularity question: a representation can be learned from a full history while the location of the event remains unknown.

\subsection{Reconstruction, weak supervision, and aggregation}

Autoencoders are often used either to reduce dimensionality or to score deviations from reconstructed normal behaviour. In the verified recent set, Pamir et al. use a gated convolutional autoencoder as a feature extractor inside a supervised, cost-sensitive LSTM pipeline rather than treating reconstruction error alone as the final theft decision \cite{Pamir2023EnergyOptimization}. That distinction is important. Reconstruction objectives reward faithful reproduction of common structure; they do not necessarily align with the ranking of theft-associated customers. AMP's own frozen reconstruction experiments reached the same methodological conclusion and motivated a customer-supervised formulation.

Multiple-instance learning provides a principled vocabulary for customer bags and unlabeled windows \cite{Carbonneau2018MILSurvey}. However, an MIL architecture still needs an aggregation rule. Mean, maximum, and attention pooling each produce a single vector and can discard distributional structure. AMP therefore treats weak supervision and aggregation as separate design questions: the loss remains customer-level, while the representation retains several statistics of the customer's latent-window distribution.

\subsection{Class imbalance and evaluation}

The theft class is a minority in SGCC and many operational datasets. He and Garcia review data-level, algorithm-level, and hybrid responses to imbalanced learning \cite{He2009Imbalanced}. Recent theft studies use class weighting, undersampling, SMOTE variants, or generative procedures, but those choices can change the effective training distribution and complicate comparison. Saito and Rehmsmeier show that precision--recall plots are more informative than ROC plots alone when positives are rare \cite{Saito2015PR}. AMP therefore treats ROC-AUC and Average Precision as co-primary customer-level metrics and avoids accuracy as a primary measure.
